% ---------------------------------------------------------------------
% Argomento: 3 ingressi, 2 porte AND/OR in cascata, un ingresso negato (pallino) sulla prima porta
%   fila A: X = AND(A, NOT B),  Y = OR(X, C)
%   fila B: X = OR(NOT A, B),  Y = AND(X, C)
%   fila C: X = OR(B, NOT C),  Y = AND(X, A)
% Solo tabella di verita' (le espressioni logiche non sono ancora state spiegate)
% Livello: intermedio   |   Tempo stimato: 8 min
% ---------------------------------------------------------------------
\begin{esercizio}[6]
Scrivi la tabella di verità del circuito in figura (colonne $A$, $B$, $C$, $X$, $Y$).

\begin{center}
\ABC{%
  \begin{circuitikz}[x=1cm,y=1cm,scale=1.25,transform shape]
    \node[american and port] (g1) at (0,0.3) {};
    \node[american or port] (g2) at (3,-0.25) {};
    \draw (g1.in 1) -- ++(-0.9,0) node[left] {$A$};
    \draw (g1.in 2) -- ++(-0.9,0) node[left] {$B$};
    \draw (g2.in 2) -- (g1.in 2 |- g2.in 2) -- ++(-0.9,0) node[left] {$C$};
    \draw (g1.out) -- ++(0.45,0) node[above] {$X$} |- (g2.in 1);
    \draw (g2.out) -- ++(0.7,0) node[right] {$Y$};
    \node[circle,draw,fill=white,inner sep=0pt,minimum size=1.6mm,anchor=east] at (g1.bin 2) {};
  \end{circuitikz}}{%
  \begin{circuitikz}[x=1cm,y=1cm,scale=1.25,transform shape]
    \node[american or port] (g1) at (0,0.3) {};
    \node[american and port] (g2) at (3,-0.25) {};
    \draw (g1.in 1) -- ++(-0.9,0) node[left] {$A$};
    \draw (g1.in 2) -- ++(-0.9,0) node[left] {$B$};
    \draw (g2.in 2) -- (g1.in 2 |- g2.in 2) -- ++(-0.9,0) node[left] {$C$};
    \draw (g1.out) -- ++(0.45,0) node[above] {$X$} |- (g2.in 1);
    \draw (g2.out) -- ++(0.7,0) node[right] {$Y$};
    \node[circle,draw,fill=white,inner sep=0pt,minimum size=1.6mm,anchor=east] at (g1.bin 1) {};
  \end{circuitikz}}{%
  \begin{circuitikz}[x=1cm,y=1cm,scale=1.25,transform shape]
    \node[american or port] (g1) at (0,-0.3) {};
    \node[american and port] (g2) at (3,0.25) {};
    \draw (g1.in 1) -- ++(-0.9,0) node[left] {$B$};
    \draw (g1.in 2) -- ++(-0.9,0) node[left] {$C$};
    \draw (g2.in 1) -- (g1.in 1 |- g2.in 1) -- ++(-0.9,0) node[left] {$A$};
    \draw (g1.out) -- ++(0.45,0) node[below] {$X$} |- (g2.in 2);
    \draw (g2.out) -- ++(0.7,0) node[right] {$Y$};
    \node[circle,draw,fill=white,inner sep=0pt,minimum size=1.6mm,anchor=east] at (g1.bin 2) {};
  \end{circuitikz}}
\end{center}

\begin{solution}
\ABC{%
  {\renewcommand{\arraystretch}{1.05}\small
  \begin{tabular}{ccc|c|c|c}
    $A$ & $B$ & $C$ & NOT $B$ & $X$ & $Y$\\\hline
    0 & 0 & 0 & 1 & 0 & 0\\
    0 & 0 & 1 & 1 & 0 & 1\\
    0 & 1 & 0 & 0 & 0 & 0\\
    0 & 1 & 1 & 0 & 0 & 1\\
    1 & 0 & 0 & 1 & 1 & 1\\
    1 & 0 & 1 & 1 & 1 & 1\\
    1 & 1 & 0 & 0 & 0 & 0\\
    1 & 1 & 1 & 0 & 0 & 1\\
  \end{tabular}}}{%
  {\renewcommand{\arraystretch}{1.05}\small
  \begin{tabular}{ccc|c|c|c}
    $A$ & $B$ & $C$ & NOT $A$ & $X$ & $Y$\\\hline
    0 & 0 & 0 & 1 & 1 & 0\\
    0 & 0 & 1 & 1 & 1 & 1\\
    0 & 1 & 0 & 1 & 1 & 0\\
    0 & 1 & 1 & 1 & 1 & 1\\
    1 & 0 & 0 & 0 & 0 & 0\\
    1 & 0 & 1 & 0 & 0 & 0\\
    1 & 1 & 0 & 0 & 1 & 0\\
    1 & 1 & 1 & 0 & 1 & 1\\
  \end{tabular}}}{%
  {\renewcommand{\arraystretch}{1.05}\small
  \begin{tabular}{ccc|c|c|c}
    $A$ & $B$ & $C$ & NOT $C$ & $X$ & $Y$\\\hline
    0 & 0 & 0 & 1 & 1 & 0\\
    0 & 0 & 1 & 0 & 0 & 0\\
    0 & 1 & 0 & 1 & 1 & 0\\
    0 & 1 & 1 & 0 & 1 & 0\\
    1 & 0 & 0 & 1 & 1 & 1\\
    1 & 0 & 1 & 0 & 0 & 0\\
    1 & 1 & 0 & 1 & 1 & 1\\
    1 & 1 & 1 & 0 & 1 & 1\\
  \end{tabular}}}\\[4pt]
  Il pallino all'ingresso di una porta nega quel segnale prima della porta; il pallino all'uscita
  (NAND, NOR) nega il risultato della porta.
  \criterioABC
  {Le 8 combinazioni di $A$, $B$, $C$ tutte presenti: 1 pt; colonna $X$ (attenzione a NOT $B$): 2 pt; colonna $Y$: 3 pt ($-0{,}5$ per ogni riga sbagliata; un errore su $X$ non si ripenalizza su $Y$ se $Y$ è coerente). Colonne attese da 000 a 111: $X$ = 00001100, $Y$ = 01011101.}
  {Le 8 combinazioni di $A$, $B$, $C$ tutte presenti: 1 pt; colonna $X$ (attenzione a NOT $A$): 2 pt; colonna $Y$: 3 pt ($-0{,}5$ per ogni riga sbagliata; un errore su $X$ non si ripenalizza su $Y$ se $Y$ è coerente). Colonne attese da 000 a 111: $X$ = 11110011, $Y$ = 01010001.}
  {Le 8 combinazioni di $A$, $B$, $C$ tutte presenti: 1 pt; colonna $X$ (attenzione a NOT $C$): 2 pt; colonna $Y$: 3 pt ($-0{,}5$ per ogni riga sbagliata; un errore su $X$ non si ripenalizza su $Y$ se $Y$ è coerente). Colonne attese da 000 a 111: $X$ = 10111011, $Y$ = 00001011.}
\end{solution}
\end{esercizio}
